\scanpage{180}
% archive sheet number: 179
celà résulte encore du fait que les fonctions
régulières sur $L(W)$ s'appliquent \uline{sur} les fonctions
régulières sur $G$,

\diagram{\resizebox{\linewidth}{!}{\begin{tikzcd}[ampersand replacement=\&,column sep=0.3em,row sep=1.0em]
\text{$K$-catégorie} \& \Longleftrightarrow \&
K\text{-}\overset{\inserted{\text{hyperalgèbre}}}{\deleted{\text{algèbre}}}\,\underline U
\arrow[r,Leftrightarrow] \& \substack{\text{hyperalgèbre commu-}\\\text{tative }A,\ \text{réun.}\\\text{filtrante }\ldotp\ldotp\ldotp\ldotp}\\
\& \substack{\text{\uline{cas particuliers}}\\[3pt]\text{\uline{Groupe pseudo-algébrique}}\\
(\text{cas où }G(\underline U)\ \deleted{\text{est}}\ \inserted{\text{total}}\ \text{dans }\underline U,\\
\text{ou encore }\underline A\ \text{soit }\unsure{p|\rho}\text{-radical, i.e.}\\
\text{l'intersection des noyaux des caractères}\\\text{de }A\text{ est réduite à }0)}
\arrow[u,Rightarrow] \& \& \\
\& \substack{\text{Groupe algébrique}\\(\text{si }\deleted{\unsure{\text{des}|\text{de}}}\ \text{plus }A\text{ est une }k\text{-}\\\text{algèbre de type fini})} \arrow[u,Rightarrow] \& \&
\end{tikzcd}}}
\margin{Si $A$ est une $k$-algèbre de type fini, \inserted{sur $k$ alg. clos,} pour qu'il s'agisse d'un groupe algébrique, (il faut et) il suffit que $A$ n'ait pas d'\unsure{élémts|éléments} nilpotents $\ne0$ ($k$ alg\textsuperscript{t.} clos est-il nécessaire ?) Ceci est-il automatiquement vérifié en caractéristique $0$ ???}
